\subsection{Evaluation Criteria}

The evaluation criteria are designed to capture key aspects of reasoning in ill-defined tasks, including problem structuring, solution generation and reflective reasoning (see Table \ref{table:evaluation_criteria}). Each criterion is rated on a 5-point Likert scale:

\begin{itemize}
    \item 5 - excellent, no major weaknesses
    \item 4 - good, minor issues
    \item 3 - acceptable, noticeable gaps
    \item 2 - poor, major issues
    \item 1 - very poor or missing
\end{itemize}

\begin{table}[h]
\centering
\caption{Evaluation criteria for ill-defined task assessment}
\begin{tabularx}{\textwidth}{l | l X}
\midrule
\textbf{Category} & \textbf{Indicator} & \textbf{Description} \\
\midrule

\multirow{3}{*}{Problem Structuring} 
& Framing 
& Identification of key elements, goals, and assumptions \\

\cmidrule{2-3}

& Decomposition 
& Degree to which the problem is structured into coherent subproblems \\

\cmidrule{2-3}

& Uncertainty Handling 
& Ability to address ambiguity, missing information, and assumptions \\

\midrule

\multirow{3}{*}{Solution Quality} 
& Coherence 
& Logical consistency and clarity of reasoning steps \\

\cmidrule{2-3}

& Justification 
& Strength of argumentation and consideration of trade-offs \\

\cmidrule{2-3}

& Knowledge Integration 
& Use of relevant domain knowledge and contextual information \\

\midrule

\multirow{2}{*}{Reflective Reasoning} 
& Diversity 
& Exploration of alternative solutions and perspectives \\

\cmidrule{2-3}

& Metacognition 
& Reflection on reasoning process and awareness of limitations \\

\midrule
\end{tabularx}
\label{table:evaluation_criteria}
\end{table}